% ============================================================
%  基于监督学习的国标麻将AI优化：特征工程、网络结构与推理后处理
%  ICML 2026 template (single submission, Chinese)
%  编译方式：xelatex -> bibtex -> xelatex -> xelatex
%  或：xelatex mahjong_sl_report.tex (两次以生成目录/交叉引用)
% ============================================================

\documentclass{article}

% ---------- ICML-style packages ----------
\usepackage[utf8]{inputenc}
\usepackage[T1]{fontenc}
\usepackage{xeCJK}             % XeLaTeX 中文支持（比 CJKutf8 更稳定）
\setCJKmainfont{SimSun}        % 宋体；可换 FandolSong / NotoSerifCJKsc
\setCJKsansfont{SimHei}
\setCJKmonofont{FangSong}
\usepackage{times}
\usepackage{microtype}
\usepackage{graphicx}
\usepackage{booktabs}
\usepackage{amsmath,amssymb}
\usepackage{algorithm}
\usepackage{algorithmic}
\usepackage[colorlinks=true,linkcolor=blue,citecolor=blue,urlcolor=blue]{hyperref}
\usepackage{url}
\usepackage{xcolor}
\usepackage{geometry}
\usepackage{enumitem}
\usepackage{multicol}
\usepackage{float}
\usepackage{caption}
\usepackage{subcaption}
\usepackage{tikz}
\usetikzlibrary{shapes.geometric,arrows.meta,positioning,fit,backgrounds}
\usepackage{pgfplots}
\pgfplotsset{compat=1.18}

% ---------- Page layout (ICML-like, A4) ----------
\geometry{
  a4paper,
  top=2.4cm, bottom=2.4cm,
  left=1.8cm, right=1.8cm
}
\setlength{\columnsep}{0.75cm}

% ---------- Handy macros ----------
\newcommand{\bx}{\mathbf{x}}
\newcommand{\calL}{\mathcal{L}}
\newcommand{\calB}{\mathcal{B}}
\newcommand{\RR}{\mathbb{R}}

% ============================================================
\title{%
  \textbf{基于监督学习的国标麻将AI优化}\\
  \large 特征工程、网络结构与推理后处理
}

\author{%
  匿名作者 \\
  北京大学 \\
  \texttt{Botzone 平台}
}
\date{2026 年 6 月}

% ============================================================
\begin{document}
\maketitle
\thispagestyle{empty}

% -------- Abstract --------
\begin{abstract}
国标麻将（Chinese Standard Mahjong，CSM）是一个具有高度不确定性、部分可观测性和
多智能体竞争特性的复杂棋牌游戏。本文系统总结了在 Botzone 平台上构建和优化 CSM
监督学习（Supervised Learning，SL）Bot 的完整过程。
我们从仅 $6$ 通道的基础特征出发，逐步扩展至 $66$ 通道的丰富特征表示，
引入基于残差网络（ResNet）的分头（multi-head）神经网络，
并设计了结合危险度估计、向听数感知和特殊牌型检测的推理后处理框架。
同时，我们探索了赢家数据过滤、番数加权采样等训练技巧，
并在华为 Ascend NPU 上完成实际训练，解决了若干平台特有的数值稳定性问题。
实验观察表明，特征扩展与规则后处理对 Bot 行为影响显著，
并发现多个有趣现象——碰/杠行为消失、点炮率过高、胡牌番数偏低——
本文分析了其根因并给出系统性解法。
\end{abstract}

\vspace{-4pt}
\begin{multicols}{2}

% ============================================================
\section{引言}
\label{sec:intro}
% ============================================================

国标麻将是中国最流行的棋牌游戏之一，规则复杂性远超日本麻将：
共 81 种番型、144 张牌、4 名玩家同场竞技，
并存在吃、碰、杠等副露操作和最低 8 番的胡牌门槛。
与完全信息博弈（如围棋）不同，麻将属于部分可观测马尔可夫决策过程（POMDP），
传统搜索方法难以直接应用~\cite{silver2016mastering}。

近年来强化学习（RL）和监督学习（SL）在麻将 AI 领域取得显著进展。
Suphx~\cite{li2020suphx} 在日本麻将上超越人类专业选手；
IJCAI 2020 竞赛~\cite{ijcai2020} 为国标麻将提供了标准数据集与评测环境。
本工作基于 IJCAI 2020 数据集，聚焦于 SL Bot 的系统性优化。

\paragraph{主要贡献}
\begin{enumerate}[leftmargin=*,topsep=2pt,itemsep=1pt,parsep=0pt]
  \item 将观测通道数从 6 扩展至 66，涵盖完整可见信息；
  \item 设计基于 ResNet 的分头模型，引入动作类型辅助预测头；
  \item 提出向听数感知的危险度惩罚与现物奖励后处理框架；
  \item 引入七对/清一色专用特征通道与番数加权训练；
  \item 系统分析多个行为异常现象的根因，具有较强的工程参考价值。
\end{enumerate}

% ============================================================
\section{背景与数据集}
\label{sec:background}
% ============================================================

\subsection{国标麻将规则摘要}

国标麻将使用 144 张牌（万/条/筒各 1--9 各 4 张，外加字牌和花牌），
4 名玩家轮流摸牌打牌。
动作空间包括弃牌（34 种）、吃（63 种组合）、碰（34 种）、
杠（明杠/暗杠/补杠各 34 种）、胡牌和过，共 \textbf{235} 个合法动作。
胡牌需满足最低 8 番的门槛，鼓励玩家追求高番牌型。

\subsection{数据集描述}

训练数据来自 IJCAI 2020 冠军 Bot 的自对弈，共 \textbf{98,209} 局。
参考实现为中科大团队（比赛第 4 名）的纯 SL 方案：
状态空间 $4\times9\times145$ 通道，动作空间 235 维。

% ============================================================
\section{算法设计}
\label{sec:method}
% ============================================================

\subsection{观测特征工程}

\begin{table}[H]
\centering
\caption{66 通道特征表示（最终版本）}
\label{tab:features}
\small
\setlength{\tabcolsep}{4pt}
\begin{tabular}{clr}
\toprule
通道区间 & 含义 & 数量 \\
\midrule
0       & 门风牌（one-hot 广播）     & 1  \\
1       & 圈风牌（one-hot 广播）     & 1  \\
2--5    & 手牌计数（第 $k$ 层代表 $\geq\!k\!+\!1$ 张）& 4  \\
6--9    & 场上可见牌（弃牌+副露）    & 4  \\
10--13  & 剩余牌数（归一化）         & 4  \\
14      & 最近一张弃牌               & 1  \\
15--18  & 各玩家牌墙剩余（归一化）   & 4  \\
19--34  & 各玩家弃牌历史（$4\times4$）& 16 \\
35--50  & 各玩家副露牌（$4\times4$） & 16 \\
51      & 杠后摸牌标志               & 1  \\
52      & 最后一张牌标志             & 1  \\
53--55  & 对手危险度（3 位对手）     & 3  \\
56--59  & 弃牌位置编码（归一化序号） & 4  \\
60      & 向听数（归一化，$0\!=\!$听牌）& 1  \\
61      & 有效进张牌（归一化剩余数） & 1  \\
62      & 七对向听数（归一化）       & 1  \\
63--65  & 万/条/筒清一色向听数       & 3  \\
\midrule
        & \textbf{合计}              & \textbf{66} \\
\bottomrule
\end{tabular}
\end{table}

特征矩阵维度为 $C \times 4 \times 9$（$C=66$），
$4\times9$ 对应麻将牌的自然二维网格（4 种花色，每色 9 张）。

\paragraph{危险度估计（通道 53--55）}
对每位对手 $p$ 及候选弃牌 $t$，采用启发式评分：
\begin{equation}
d(p,t) = \alpha_0 \mathbb{1}[\text{非现物}]
        + \alpha_1 \mathbb{1}[\text{筋牌}]
        + \alpha_2 \cdot \text{suit\_pressure}(p)
\label{eq:danger}
\end{equation}
其中现物（已被对手打出的牌）风险为 0，
副露牌花色附近的牌风险更高（$\alpha_2$ 随副露数线性增加）。

\paragraph{有效进张（通道 61）}
穷举 34 种牌类测试向听数变化：
若摸入牌 $t$ 后向听数下降，则该格填入剩余张数 $\text{remaining}(t)/4$；
否则为 0。最多调用 34 次 \texttt{MahjongShanten}，
在预处理阶段性能可接受。

\paragraph{特殊番型向听（通道 62--65）}
七对向听 $= 6 - |\text{对子数}|$；
清一色向听粗估为 $13 - |\text{同花色手牌数}|$，
归一化后分别写入对应通道。

\subsection{神经网络结构}

网络采用 ResNet 风格的分头架构，如图~\ref{fig:model} 所示。

\paragraph{主干（Trunk）}
两层卷积提取局面特征：
$\text{Conv}_{3\times3}^{128} \to \text{BN} \to \text{ReLU}
 \to \text{Conv}_{3\times3}^{64} \to \text{BN} \to \text{ReLU}$。

\paragraph{残差块（ResBlock $\times$ 6）}
每块结构为：
$\text{Conv}_{3\times3}^{64} \to \text{BN} \to \text{ReLU}
 \to \text{Conv}_{3\times3}^{64} \to \text{BN} \to \oplus \to \text{ReLU}$，
其中 $\oplus$ 为残差加和，带 $1\times1$ 跳接投影层~\cite{he2016deep}。

\paragraph{分头（Multi-Head）}
Flatten $\to$ FC512 $\to$ ReLU $\to$ FC256 $\to$ ReLU 之后分叉：
\begin{itemize}[leftmargin=*,topsep=2pt,itemsep=1pt,parsep=0pt]
  \item \textbf{动作类型头}（type\_head）：FC8，预测 8 类动作类型
    （Pass/Hu/Play/Chi/Peng/Gang/AnGang/BuGang）；
  \item \textbf{动作子头}（6 个 tile\_head）：每类动作各一个 FC，
    输出具体牌位 logit。
\end{itemize}

\paragraph{Logit 合成}
\begin{equation}
\text{logit}(a) = \text{typeLogit}[\text{class}(a)] + \text{tileLogit}(a)
\end{equation}
非法动作位填 $-100.0$ 后做 softmax。选择 $-100$ 而非 $-\infty$ 是为了
避免与 label smoothing 交互时 $\log\text{softmax}(-\infty)$ 溢出。

\paragraph{训练目标}
\begin{equation}
\calL = \calL_{\text{policy}} + 0.2\,\calL_{\text{type}}
\end{equation}
$\calL_{\text{policy}}$ 为番数加权交叉熵（见第~\ref{sec:training} 节），
$\calL_{\text{type}}$ 为动作类型的交叉熵辅助损失。

\subsection{推理后处理框架}

推理时对模型原始 logit 依次施加以下调整（图~\ref{fig:postprocess}）：

\paragraph{① 向听数感知危险惩罚}
\begin{equation}
a_{\text{play}} \mathrel{-}= 2.0 \cdot d(t) \cdot
\min\!\left(1,\; \frac{\text{shanten}}{3}\right)
\label{eq:danger_pen}
\end{equation}
向听数 $\geq 3$ 时全力防守，向听数 $=0$（听牌）时不惩罚，
允许进攻摸胡。

\paragraph{② 现物（Genbutsu）奖励}
\begin{equation}
a_{\text{play}} \mathrel{+}= 0.3
\quad \text{若 } t \in \bigcup_{p=1}^{3}\text{discards}(p)
\label{eq:genbutsu}
\end{equation}
弱奖励（$+0.3$）确保模型有明显偏好时仍能覆盖规则建议，
避免完全服从现物而放弃牌效。

\paragraph{③ 听牌进攻}
若 $\text{shanten}=0$，则 $\text{Pass} \mathrel{-}= 1.0$，
抑制消极等待行为。

\paragraph{④ 七对子保护}
若七对向听 $\leq$ 普通向听且 $\leq 2$，
对子牌打出得 $-0.5$ 惩罚，引导 Bot 保留对子追七对。

\paragraph{⑤ 轻碰/杠惩罚}
吃 $-0.45$，碰 $-0.35$，杠 $-0.30$，补杠 $-0.25$，
平衡副露积极性（过强惩罚会导致碰/杠行为消失，见第~\ref{sec:phenomenon} 节）。

% ============================================================
\section{训练流程优化}
\label{sec:training}
% ============================================================

\subsection{赢家数据过滤}

解析每局 \texttt{Score} 行，仅保留得分最高玩家的状态-动作对。
直觉上胡牌玩家的决策更具参考价值；
但完全过滤可能导致防守样本稀缺，是值得深入研究的权衡点。

\subsection{番数加权采样}

用赢家得分代理番数，计算样本权重：
\begin{equation}
w_i = \log_2\!\left(\frac{\max(\text{score}_i,\,8)}{8}\right) + 1
\end{equation}
8 分对应权重 1.0，128 分对应权重 4.0。加权策略损失为：
\begin{equation}
\calL_{\text{policy}}
= \frac{1}{|\calB|}\sum_{i\in\calB} \bar{w}_i \cdot \text{CE}(f(s_i),a_i)
\label{eq:weighted_loss}
\end{equation}
其中 $\bar{w}_i = w_i / \bar{w}$ 为 batch 内归一化权重。

\subsection{延迟 LRU 加载}

特征扩展后（6$\to$66 通道），全量加载约需 30 GB 内存。
改用 \texttt{OrderedDict} 实现 LRU 缓存（默认 256 个 \texttt{.npz} 文件），
内存峰值降至 $\sim$4 GB。

\subsection{多进程预处理}

使用 \texttt{multiprocessing.Pool.imap()} 并行处理，
实测加速比约 $N_{\text{core}}-1$ 倍。
工程注意事项：连续特征须以 float16 保存（int8 会将 $[0,1]$ 截断为 0）。

\subsection{断点续训}

保存完整训练状态（model / optimizer / scheduler / epoch），
支持 \texttt{--resume} 从任意 checkpoint 继续，
在云端 NPU 环境（任务超时频繁中断）中尤为重要。

% ============================================================
\section{有趣的现象与启发}
\label{sec:phenomenon}
% ============================================================

\subsection{碰/杠行为消失}

\textbf{现象：}早期版本 Bot 完全不碰牌、不杠牌。

\textbf{根因分析：}两个耦合因素：
① 赢家数据偏门清（清一色等番型往往门清不副露），
   导致训练集碰/杠样本极少；
② 后处理惩罚过重（原始碰 $-0.80$，杠 $-0.80$），
   彻底掩盖模型的原始碰牌偏好。

\textbf{解法：}将惩罚降至碰 $-0.35$、杠 $-0.30$，
并对持有 3 张以上时额外回调奖励，恢复正常副露行为。

\textbf{启发：}\textit{后处理强规则与模型分布之间的耦合效应非常显著}——
过强的规则惩罚可能完全覆盖模型学到的有效策略，
需在每次调参后进行行为层面的验证，而不仅仅是看 loss 或准确率。

\subsection{频繁点炮问题}

\textbf{现象：}Bot 以较高频率打出对手需要的胡牌牌（点炮）。

\textbf{根因分析：}纯 SL 模型从胡牌玩家视角学习，
缺乏"哪张牌对对手危险"的显式感知；
PACK 通道虽已编码对手副露，但模型尚未建立可靠的危险度因果关联。

\textbf{解法：}
引入危险度特征通道（通道 53--55）并结合式~(\ref{eq:danger_pen}) 的向听感知衰减——
越接近听牌越少防守，实现进攻与防守的动态平衡。
弱奖励安全牌（式~\ref{eq:genbutsu}，$+0.3$）比强惩罚危险牌更稳健，
因为前者不会覆盖模型在牌效上的合理判断。

\subsection{平均番数偏低}

\textbf{现象：}Bot 胡牌平均番数较低，七对、清一色等高番牌型极少出现。

\textbf{根因分析：}
① 高番局在训练数据中占比小（$<5\%$）；
② 模型缺乏对七对、清一色等特殊番型向听进度的感知；
③ 未加权的交叉熵损失对低番高频样本过度拟合。

\textbf{解法（已实现，重训后生效）：}
七对/清一色向听数特征通道（通道 62--65）+
番数加权损失（式~\ref{eq:weighted_loss}）+
七对子后处理保护（$-0.5$ 惩罚打对子）。

\subsection{NPU 平台数值稳定性问题}

在华为 Ascend NPU 上训练时发现若干平台特有问题，
对跨平台部署具有较强参考价值（表~\ref{tab:npu}）。

\begin{table}[H]
\centering
\caption{NPU 平台典型问题与解法}
\label{tab:npu}
\small
\setlength{\tabcolsep}{4pt}
\begin{tabular}{p{2.2cm}p{2.0cm}p{2.2cm}}
\toprule
问题 & 现象 & 解法 \\
\midrule
\texttt{non\_blocking=True} &
  error 507001 Segfault &
  NPU 强制关闭 \\
masked\_fill 用 finfo.min &
  loss 爆炸至 $10^7$ &
  改用 $-100.0$ \\
int8 保存连续特征 &
  归一化值截断为 0 &
  改用 float16 \\
label\_smooth $\times$ 大负 mask &
  log-softmax($-\infty$)溢出 &
  mask 值 $\geq -100$ \\
\bottomrule
\end{tabular}
\end{table}

\textbf{启发：}跨硬件平台迁移时，数值稳定性问题往往比 API 兼容性更难调试；
建议对关键张量的统计量（min/max/NaN 比例）进行实时监控。

% ============================================================
\section{开放式问题与研究方向}
\label{sec:open}
% ============================================================

\subsection{SL 上限与数据多样性}

纯 SL 方法的上限受制于训练数据的策略分布。
冠军 Bot 存在风格偏差（偏门清），使赢家过滤后副露样本稀缺。
\textbf{研究方向：}多样化数据源混合（不同风格 Bot 的对局），
或设计基于结果的重采样策略——在保留高质量赢家样本的同时，
补充足够的防守样本以平衡策略分布。

\subsection{从 SL 到 RL 的迁移}

SL 模型是 RL 的良好初始化策略。
参考 AlphaGo~\cite{silver2016mastering} 的经验，
SL 策略网络 $\to$ RL 自博弈可进一步提升。
\textbf{挑战：}CSM 的稀疏奖励（仅局末才知道得分）$+$ 4 人博弈的信用分配
$+$ 番数计算的规则复杂性，使奖励函数设计本身就是核心研究课题。

一个有趣的研究问题：如何设计稠密中间奖励？
候选方案：向听数变化奖励、危险度降低奖励、胡牌番数奖励。
这些中间奖励是否会产生"奖励黑客"行为（例如学会无限循环保持听牌状态）？

\subsection{隐藏信息的贝叶斯推断}

当前特征仅使用可见信息。对手手牌的贝叶斯推断是提升防守能力的重要方向：
\begin{equation}
P(\text{hand}_j \mid \text{obs}_{1:t})
\propto P(\text{obs}_{1:t} \mid \text{hand}_j)\, P(\text{hand}_j)
\end{equation}
这本质上是基于不完整观测的手牌推断问题，
可用变分推断或粒子滤波建模，是麻将 AI 区别于完全信息博弈的核心挑战。

\subsection{后处理的参数化与可学习化}

当前后处理系统包含约 10 个人工调参的超参数，
存在过拟合当前模型偏差的风险。
\textbf{研究方向：}
用轻量残差网络替代规则后处理（"可微规则"），
或通过少量对弈数据学习后处理层参数，实现端到端优化。

\subsection{番数预测辅助任务}

可进一步引入\textbf{番数预测头}（预测最终胡牌番数期望），
作为价值估计辅助信号，为危险惩罚的衰减系数提供数据驱动的依据，
替代当前式~(\ref{eq:danger_pen}) 中的硬编码阈值 3。

% ============================================================
\section{结论}
\label{sec:conclusion}
% ============================================================

本文系统总结了国标麻将 SL Bot 从基础实现到工程化优化的完整过程，
主要贡献包括：66 通道精细特征工程、ResNet 分头模型、
向听数感知的推理后处理框架、赢家过滤与番数加权训练，
以及多项 Ascend NPU 工程适配。

实践表明，\textbf{特征工程与后处理规则对 SL Bot 行为的影响十分显著，
有时甚至超过模型容量本身}。
多个有趣的行为异常（碰牌消失、频繁点炮、高番稀少）
都能追溯到数据分布偏差或后处理参数不当，
提示在 AI 棋牌游戏研究中需对数据质量和推理部署给予同等重视。

未来工作的主要方向是从 SL 过渡到 RL 自博弈，
引入隐藏信息推断模块，以及将规则后处理参数化为可学习组件。

\end{multicols}

% ============================================================
% Figures (full-width, placed after two-column body)
% ============================================================

\begin{figure}[h]
\centering
\begin{tikzpicture}[
  >=Latex,
  box/.style={
    rectangle, draw=black!70, rounded corners=3pt,
    minimum width=3.2cm, minimum height=0.65cm,
    text centered, font=\small
  },
  sbox/.style={box, minimum width=2.4cm, fill=red!10},
  node distance=0.45cm and 0.5cm
]
% Main trunk
\node[box, fill=blue!12] (inp) {输入 $66\times4\times9$};
\node[box, fill=blue!12, right=of inp] (c1)  {Conv$_{3\times3}^{128}$ + BN + ReLU};
\node[box, fill=blue!12, right=of c1]  (c2)  {Conv$_{3\times3}^{64}$ + BN + ReLU};
\node[box, fill=green!15, right=of c2] (res) {ResBlock $\times$ 6};
\node[box, fill=orange!15, right=of res] (fc) {FC512 $\to$ FC256};

% Heads
\node[sbox, below left=0.8cm and 0.8cm of fc]  (th)  {type\_head (FC8)};
\node[sbox, below right=0.8cm and 0.5cm of fc] (dh)  {6 tile sub-heads};
\node[box, fill=gray!20, below=1.8cm of fc]     (out) {Merge + Mask $\to$ 235-dim logit};

% Arrows
\draw[->] (inp) -- (c1);
\draw[->] (c1)  -- (c2);
\draw[->] (c2)  -- (res);
\draw[->] (res) -- (fc);
\draw[->] (fc)  -- (th);
\draw[->] (fc)  -- (dh);
\draw[->] (th)  -- (out);
\draw[->] (dh)  -- (out);
\end{tikzpicture}
\caption{分头 ResNet 模型结构。主干提取局面特征，
type\_head 提供动作类型辅助监督，6 个 tile sub-head 预测具体动作 logit，
最终合并后施加合法动作掩码输出 235 维策略向量。}
\label{fig:model}
\end{figure}

\begin{figure}[h]
\centering
\begin{tikzpicture}[
  >=Latex,
  step/.style={rectangle, draw=black!70, rounded corners=3pt,
    fill=yellow!10, minimum width=11cm, minimum height=0.65cm,
    text centered, font=\small},
  node distance=0.38cm
]
\node[step] (raw)  {原始模型 logit（235 维）};
\node[step, below=of raw]  (s1) {① 向听感知危险惩罚（式~\ref{eq:danger_pen}）：$a_{\text{play}} \mathrel{-}= 2.0 \cdot d(t) \cdot \min(1, \text{shanten}/3)$};
\node[step, below=of s1]   (s2) {② 现物奖励（式~\ref{eq:genbutsu}）：安全牌 $+0.3$};
\node[step, below=of s2]   (s3) {③ 听牌进攻：$\text{shanten}=0$ 时 Pass $-1.0$};
\node[step, below=of s3]   (s4) {④ 七对子保护：对子牌 $-0.5$（七对向听 $\leq 2$ 时）};
\node[step, below=of s4]   (s5) {⑤ 轻碰/杠惩罚：吃$-0.45$，碰$-0.35$，杠$-0.30$，补杠$-0.25$};
\node[step, below=of s5, fill=green!10]   (end) {非法动作填 $-100.0$，取 argmax 输出动作};

\draw[->] (raw) -- (s1);
\draw[->] (s1)  -- (s2);
\draw[->] (s2)  -- (s3);
\draw[->] (s3)  -- (s4);
\draw[->] (s4)  -- (s5);
\draw[->] (s5)  -- (end);
\end{tikzpicture}
\caption{推理后处理流水线。每步对 logit 向量进行原地调整，
最终取 argmax 输出合法动作。各步调整量均经过人工验证，
确保模型偏好强时规则调整不覆盖模型判断。}
\label{fig:postprocess}
\end{figure}

% ============================================================
\bibliographystyle{plain}
\begin{thebibliography}{10}

\bibitem{li2020suphx}
Junjie Li, Shuai Ma, Malong Yang, Qiaoyuan Mo, Ning Ding, Qiang Yan, and Yuandong Tian.
\newblock Suphx: Mastering mahjong with deep reinforcement learning.
\newblock \textit{arXiv preprint arXiv:2003.13590}, 2020.

\bibitem{silver2016mastering}
David Silver, Aja Huang, Chris~J. Maddison, Arthur Guez, Laurent Sifre,
  George van den Driessche, Julian Schrittwieser, Ioannis Antonoglou,
  Veda Panneershelvam, Marc Lanctot, Sander Dieleman, Dominik Grewe,
  John Nham, Nal Kalchbrenner, Ilya Sutskever, Timothy Lillicrap, Madeleine Leach,
  Koray Kavukcuoglu, Thore Graepel, and Demis Hassabis.
\newblock Mastering the game of go with deep neural networks and tree search.
\newblock \textit{Nature}, 529(7587):484--489, 2016.

\bibitem{silver2017mastering}
David Silver, Julian Schrittwieser, Karen Simonyan, Ioannis Antonoglou,
  Aja Huang, Arthur Guez, Thomas Hubert, Lucas Baker, Matthew Lai,
  Adrian Bolton, Yutian Chen, Timothy Lillicrap, Fan Hui, Laurent Sifre,
  George van den Driessche, Thore Graepel, and Demis Hassabis.
\newblock Mastering the game of go without human knowledge.
\newblock \textit{Nature}, 550(7676):354--359, 2017.

\bibitem{he2016deep}
Kaiming He, Xiangyu Zhang, Shaoqing Ren, and Jian Sun.
\newblock Deep residual learning for image recognition.
\newblock In \textit{Proceedings of the IEEE Conference on Computer Vision and
  Pattern Recognition (CVPR)}, pages 770--778, 2016.

\bibitem{moravvcik2017deepstack}
Matej Mora{\v{c}}{\'\i}k, Martin Schmid, Neil Burch, Viliam Lis{\'y},
  Dustin Morrill, Nolan Bard, Trevor Davis, Kevin Waugh, Marc Lanctot,
  and Michael Bowling.
\newblock Deepstack: Expert-level artificial intelligence in heads-up no-limit poker.
\newblock \textit{Science}, 356(6337):508--513, 2017.

\bibitem{ijcai2020}
{IJCAI-2020 Mahjong AI Competition}.
\newblock Botzone platform: Chinese standard mahjong.
\newblock \url{http://www.botzone.org.cn/game/Chinese-Standard-Mahjong}, 2020.

\bibitem{mnih2015human}
Volodymyr Mnih, Koray Kavukcuoglu, David Silver, Andrei~A. Rusu, Joel Veness,
  Marc~G. Bellemare, Alex Graves, Martin Riedmiller, Andreas~K. Fidjeland,
  Georg Ostrovski, Stig Petersen, Charles Beattie, Amir Sadik, Ioannis Antonoglou,
  Helen King, Dharshan Kumaran, Daan Wierstra, Shane Legg, and Demis Hassabis.
\newblock Human-level control through deep reinforcement learning.
\newblock \textit{Nature}, 518(7540):529--533, 2015.

\bibitem{brown2019superhuman}
Noam Brown and Tuomas Sandholm.
\newblock Superhuman {AI} for multiplayer poker.
\newblock \textit{Science}, 365(6456):885--890, 2019.

\end{thebibliography}

\end{document}
